\section{Execute Commands 与 Write Stdin}

\subsection{Execute Commands}

Execute Commands 是一次性的命令调用：启动进程、运行一段、获得输出。它同时兼容"执行模型生成的代码脚本"——模型可以先生成一段代码，再通过 Execute Commands 运行。

\begin{knowledgebox}{"Patch is All You Need"}
讲者引用了业界流行的说法：将应用程序变成命令行可交互的形式后，所有操作都可以通过 Patch（编辑）和 Command（执行）两个原语完成。
\end{knowledgebox}

\subsection{Write Stdin}

Write Stdin 用于向\textbf{正在运行}的程序持续发送输入——即与交互式 CLI（如 iPython、Node REPL、MySQL、GDB 等）进行对话，发送输入并获取增量输出。

\begin{importantbox}{Execute Commands vs Write Stdin}
\begin{itemize}
    \item \textbf{Execute Commands}：一次性调用。启动$\rightarrow$运行$\rightarrow$获取输出$\rightarrow$结束。
    \item \textbf{Write Stdin}：持续交互。向已运行的进程发送新指令，获取增量响应。
\end{itemize}
\end{importantbox}

\subsection{本章小结}

Execute Commands 覆盖了绝大多数"运行某段代码/命令"的需求，Write Stdin 则补充了交互式会话的场景，两者共同构成了 Agent 执行能力的基础。


